\subsection{A constructive positive family of rectangular Kronecker rays}

\begin{proposition}[P-F59: stable adjoint rays with small positive forms and large periods]
\label{prop:rectangularadjoint}
Let $r\ge2$ and $d\ge A\ge1$ be integers, and put
\[
 \Lambda=(rd-A,A),\qquad R=(d^r),\qquad
 k_{r,A,d}(s)=g(s\Lambda,sR,sR)\quad(s\ge0).
                                                               \tag{P-F59a}
\]
The value at zero is one. These are actual ordinary rectangular
Kronecker stretching functions. They satisfy
\[
 k_{r,A,d}(s)
       =[t^{As}]\prod_{j=2}^{r}(1-t^j)^{-1}.
                                                               \tag{P-F59b}
\]
For $j=2,\ldots,r$, set
\[
 c_j=\gcd(A,j),\qquad L_j=A/c_j,\qquad b_j=j/c_j.
\]
Their generating series has the explicit positive form
\[
 \sum_{s\ge0}k_{r,A,d}(s)T^s
   =\frac{H_{r,A}(T)}{\prod_{j=2}^{r}(1-T^{b_j})},
 \qquad
 H_{r,A}(T)=
 \sum_{\substack{0\le u_j<L_j\ (2\le j\le r)\\
                         A\mid\sum_{j=2}^{r}ju_j}}
       T^{(\sum_{j=2}^{r}ju_j)/A}.
                                                               \tag{P-F59c}
\]
In particular, $H_{r,A}\in\mathbb Z_{\ge0}[T]$, $H_{r,A}(0)=1$,
all denominator weights are at most $r$, and
\[
 \deg H_{r,A}<\sum_{j=2}^{r}b_j
                      \le\frac{r(r+1)}2-1,
 \qquad
 [T^v]H_{r,A}(T)\le A^{r-1}.
                                                               \tag{P-F59d}
\]
For the literal constituent indices used throughout this paper, let
\[
 C_r=(r-1)r+128(r-1)^3r^3.
\]
Then
\[
 s(k_{r,A,d})\le p(k_{r,A,d})
                       \le\left\lceil C_r/A\right\rceil.
                                                               \tag{P-F59e}
\]
For $r=2$ both indices are in fact zero.
The positive form in (P-F59c), and any individual value in
(P-F59b), can be constructed or computed in time polynomial in
the binary input length together with $r$; no enumeration of the
$A^{r-1}$ potential residue vectors is required.

When $A=1$, the minimum quasi-polynomial period is exactly
\[
       \operatorname{lcm}(2,3,\ldots,r)\ge2^{r-1}.
                                                               \tag{P-F59f}
\]
Thus this actual Kronecker family has large minimum periods
alongside bounded denominator weights and polynomial strict-index
bounds. No canonical h.s.o.p. statement for the Kronecker
branching ring is inferred from equality of dimensions.
\end{proposition}

\begin{proof}
Manivel, \emph{On rectangular Kronecker coefficients},
Theorem~1, printed p.~1, states
\[
 g((rd'-|\rho|,\rho),(d'^r),(d'^r))
                   =\dim S_{\rho}(\mathfrak{sl}_r)^{GL_r}
       \qquad(d'\ge|\rho|).
\]
Here $GL_r$ acts by conjugation, and its center acts trivially,
so its invariants are the $SL_r$ invariants:
\url{https://arxiv.org/pdf/0907.3351}.
Apply the theorem with $d'=sd$ and $\rho=(sA)$ for every
$s\ge1$. Its inequality holds because $d\ge A$; also
$rd-A\ge A$, so (P-F59a) uses partitions. It gives
\[
       k_{r,A,d}(s)
          =\dim\operatorname{Sym}^{sA}
                           (\mathfrak{sl}_r)^{GL_r}.
                                                               \tag{P-F59g}
\]
This is a full stretching identity, with $r,A,d$ fixed.
It is not an inequality or a limit in a dimension that changes
with $s$.

For clarity, the invariant-ring calculation in this case is
elementary. Restrict a conjugation-invariant polynomial on
traceless matrices to diagonal matrices with coordinates
$z_1,\ldots,z_r$ and $\sum z_i=0$. Restriction is injective:
matrices with distinct eigenvalues form a nonempty Zariski-open
subset and are conjugate to diagonal matrices. The restriction
is symmetric. Every symmetric polynomial on the hyperplane is
the restriction of a symmetric polynomial on the full diagonal
space: lift it arbitrarily and average its lift over the finite
permutation group. The fundamental symmetric-polynomial
elimination then expresses it in the elementary symmetric
functions $e_1,\ldots,e_r$, and $e_1=0$ on the hyperplane.
Each $e_j$ extends to the corresponding coefficient of the
characteristic polynomial of a matrix. Finally $e_2,\ldots,e_r$
are algebraically independent on this hyperplane: every monic
degree-$r$ polynomial with vanishing coefficient of degree
$r-1$ has complex roots whose sum is zero. Consequently
\[
 \mathbb C[\mathfrak{sl}_r]^{GL_r}
                  =\mathbb C[e_2,e_3,\ldots,e_r],
                    \qquad \deg e_j=j.
                                                               \tag{P-F59h}
\]
The trace pairing identifies $\mathfrak{sl}_r$ with its dual
equivariantly, so (P-F59g) and (P-F59h) prove (P-F59b), also
at $s=0$. Manivel's Corollary~2(1), printed p.~6, records the
same restricted-partition count.

To prove the positive form directly, write every exponent vector
of a monomial in the polynomial ring (P-F59h) uniquely as
\[
                    v_j=u_j+L_jw_j,
              \qquad 0\le u_j<L_j,\quad w_j\ge0.
\]
Since $jL_j=Ab_j$, its weight is a multiple of $A$ exactly
when $\sum ju_j$ is. After division of the grading by $A$,
the monomials with fixed such $u$ have generating series
$T^{\sum ju_j/A}/\prod_j(1-T^{b_j})$.
Summing over $u$ proves (P-F59c). The zero vector is the only
one contributing to its constant term. Also
\[
       \frac{\sum ju_j}{A}
             <\sum\frac{jL_j}{A}=\sum b_j,
\]
and the number of available vectors is at most
$\prod L_j\le A^{r-1}$. This proves every assertion in
(P-F59d), including the bit-size bound for the numerator.

Here is the strict-index argument, with the change of variable
made explicit. Put
$u_r(v)=[t^v]\prod_{j=2}^{r}(1-t^j)^{-1}$.
Partial fractions at roots of unity, with no polynomial part,
show that $u_r$ is a global quasi-polynomial. P-F11 applied to
its positive form with numerator one gives
$p(u_r)\le128(r-1)^3r^3\le C_r$.
Thus every constituent $q$ of $u_r$ has nonnegative coefficients
after translation by $C_r$; further translation by a nonnegative
number preserves this property. A period for $u_r$ is also a
period for $u_r(As)$. Every constituent of the latter is of the
form $q(AX)$ for one of the former constituents. Set
$P=\lceil C_r/A\rceil$. Since $AP-C_r\ge0$, the polynomial
\[
             q(A(X+P))=q\bigl(C_r+AX+(AP-C_r)\bigr)
\]
has nonnegative coefficients. This proves the upper bound for
$p$ in (P-F59e). A nonzero polynomial with nonnegative
coefficients is positive at every positive real argument,
which proves the stated bound for literal saturation as well.
When $r=2$, the function is the indicator of $2\mid As$;
its constituents are constants zero or one, so both indices
are zero.

We give the algorithm and its integer-size justification to
distinguish compact positive data from an exponentially long
residue enumeration. Write
\[
 D(t)=\prod_{j=2}^{r}(1-t^j)=\sum_{i=0}^{B}d_it^i,
      \qquad B=\sum_{j=2}^{r}j=O(r^2),\quad d_0=1.
                                                               \tag{P-F59i}
\]
The coefficients of $u_r$ satisfy the order-$B$ recurrence
$u_r(N)=-\sum_{i=1}^{B}d_iu_r(N-i)$ for $N\ge B$;
their first $B$ values are computed by the usual successive
geometric-series multiplication, using $O(rB)$ integer
operations. An order-$B$ companion matrix and binary powering
therefore compute $u_r(N)$ using a polynomial number of integer
operations in $r+\log(N+1)$.

All of these operations also have polynomial bit cost. Indeed,
$|d_i|\le2^{r-1}$ and $u_r(N)\le(N+1)^{r-1}$ by the direct
exponent-vector count. Any fundamental solution of the same
recurrence, with a single one among its first $B$ values and
zeros elsewhere, has generating series $P(t)/D(t)$, where
$\deg P<B$ and every coefficient of $P$ is at most
$2^{r-1}$ in absolute value. Its value at $N$ is consequently
at most $B2^{r-1}(N+1)^{r-1}$ in absolute value. Each entry of
the $k$th companion-matrix power is one of these fundamental
solutions at an index at most $k+B$. Hence it has
\[
             O\bigl(r+\log B+r\log(k+B+1)\bigr)
                                                               \tag{P-F59j}
\]
bits. Binary powering only uses powers with $k$ at most a
constant multiple of the desired index, and sums of $B$
products of such integers still have polynomial bit length.
This proves the claimed bit complexity for an individual value.

To construct the whole positive numerator, form
$D_A(T)=\prod_{j=2}^{r}(1-T^{b_j})$, whose degree
$B_A=\sum b_j$ is at most $B$. Compute the $B_A$ values
$u_r(As)$ for $0\le s<B_A$ by this recurrence algorithm,
and multiply their truncated series by $D_A(T)$.
Because (P-F59d) already proves $\deg H_{r,A}<B_A$, these
first $B_A$ coefficients recover the entire numerator
exactly. The bounds just proved are polynomial in
$r+\log(A+1)$; $d$ does not enter this computation.
The result is the nonnegative numerator established by the
monomial proof, even though the recurrence uses subtraction.
Neither minimum strict indices nor canonical parameter degrees
are being computed by this algorithm.

Finally, if $A=1$, the numerator is one and the series has
a pole at every primitive $j$th root for $2\le j\le r$.
A global quasi-polynomial of period $L$ has a generating
function whose poles are all $L$th roots of unity, by summing
its constituent polynomials on arithmetic progressions.
Thus each such $j$ divides every period. Conversely partial
fractions show that their least common multiple is a period.
This proves the asserted period formula. For its elementary
lower bound, write
$L=\operatorname{lcm}(1,\ldots,r)$. For every prime $p$
and every $0\le k<r$, Legendre's factorial formula gives
\[
 v_p\!\left(r\binom{r-1}{k}\right)
  =\sum_{a\ge1}\left(
     \left\lfloor\frac r{p^a}\right\rfloor
    -\left\lfloor\frac k{p^a}\right\rfloor
    -\left\lfloor\frac{r-1-k}{p^a}\right\rfloor\right)
      \le\lfloor\log_p r\rfloor=v_p(L).
\]
Each summand is zero or one, since
$r=k+(r-1-k)+1$, and all terms with $p^a>r$ vanish.
Thus $r\binom{r-1}{k}$ divides $L$. Summing the resulting
$r$ inequalities gives $r2^{r-1}\le rL$, as required.

The result concerns the positive rational form and the
strict indices of the actual Kronecker
stretching function. The dimension equality in (P-F59g)
alone does not identify its canonical branching ring with
(P-F59h) or its Veronese subring, and no such ring
identification is used.
\end{proof}
